\documentclass[11pt,a4paper]{article}

\usepackage{times,latexsym}
\usepackage{url}
\usepackage[T1]{fontenc}

%% TACL formatting package
%% Using [acceptedWithA] to generate the camera-ready final version.
\usepackage[acceptedWithA]{tacl2021v1}

%% User Packages
\usepackage{multirow}
\usepackage{enumitem}
\usepackage{booktabs}
\usepackage{graphicx}
\usepackage{amsmath, amssymb}

\begin{document}

\title{PersianIdioBench: A Comprehensive Benchmark for Persian Idiom Understanding and Generation}

\author{
  Author Name \\
  Institution Name / Country \\
  \texttt{email@domain.com}
}

\maketitle

\begin{abstract}
The interpretation of figurative language is a cornerstone of human-level reasoning, yet it remains a primary bottleneck for modern Natural Language Processing (NLP) systems. Idiomatic expressions present substantial difficulties for computational models due to their non-compositional meaning. While high-resource languages have seen progress in figurative language understanding, Persian remains critically under-resourced. This paper introduces PersianIdioBench, the first comprehensive, multi-task benchmark to rigorously evaluate Large Language Models (LLMs) on Persian idiom understanding and generation. Our benchmark, which is inspired by the CHENGYU-BENCH project for Chinese \citep{fu2025chengyubench}, features the creation of the largest curated repository of Persian idioms to date (4,352 unique expressions) alongside high-quality human annotations for evaluative connotation ($\kappa=0.68$). We formulate three evaluation tasks: (1) Connotation Classification for inherent sentiment; (2) Appropriateness in Context for semantic and pragmatic validity; and (3) Cloze-Style Idiom Selection for context-aware generation. Evaluating a suite of closed-source and open-source models reveals significant disparities. The resulting benchmark and datasets are publicly released on GitHub to serve as a standardized evaluation resource.
\end{abstract}

\section{Introduction}
The interpretation of figurative language remains a persistent challenge in Natural Language Processing (NLP). In contrast to literal expressions, where meaning is typically derived compositionally from individual components, idioms operate as holistic semantic units. Their interpretation tests the limits of compositional semantics and requires deep cultural and pragmatic knowledge, making them a fundamental benchmark for evaluating true natural language understanding across any language. In Persian, idiomatic expressions are not simply stylistic devices; they are essential to colloquial interaction, political rhetoric, and literary discourse. 
For example, the idiom \textit{hendvaneh zir-e baghal-e kasi gozashtan}, literally translating to ``putting a watermelon under someone's arm,'' has the specific, non-compositional meaning of ``flattering someone to manipulate them''. A model that processes this phrase literally will miss the intended meaning and may infer an erroneous semantic frame related to fruit or physical objects. Failure to process these expressions correctly leads to cascading errors in downstream tasks \citep{fadaee-etal-2017-data, korkontzelos-etal-2013-semeval}. Previous research indicates that excluding idioms from Persian sentiment analysis pipelines can result in an 8\% misclassification rate. For instance, a sentence might use a Persian idiom like \textit{daste goli be ab dadan} (literally: ``giving a flower bouquet to the water'') to imply that someone has made a huge mistake; a model reading this literally as a positive act of ``giving flowers'' would fail to decode the critical sentiment.
Without a dedicated understanding of these multi-word expressions (MWEs), NLP systems operating in Persian are liable to make high-confidence errors in sentiment analysis, machine translation, and dialog systems.
Despite the importance of Persian idioms, existing NLP benchmarks address them only marginally. The limited body of work evaluating idiom processing in Persian, particularly in the context of large language models, has largely focused on narrow tasks such as sentiment polarity detection \cite{dashtipour2022extending}, machine translation \cite{rezaeimanesh2025largelanguagemodelspersian}, or idiom meaning disambiguation \cite{khoshtab2025multilingualidioms}. While these studies provide useful insights, they do not assess whether models can interpret idioms within broader discourse or use them appropriately in context. As a result, available resources remain fragmented and are often limited to static lexicons or task-specific datasets. To mitigate this gap, we propose a benchmark, comparable to the CHENGYU-BENCH project for Chinese \citep{fu2025chengyubench}, comprising three tasks that extend beyond the capabilities of existing Persian idiom resources.

To mitigate this gap, this research proposes a benchmark comprising three tasks that extend beyond the capabilities of existing Persian idiom resources. These three tasks represent the core pillars of idiom processing across any language: Evaluative Connotation captures the foundational sentiment of the expression, Appropriateness in Context tests passive pragmatic comprehension within discourse, and Cloze-Style Selection evaluates active, context-aware generation. The first task, evaluative connotation, requires models to identify the sentiment of idioms. The second task, contextual appropriateness, assesses whether a candidate idiom fits naturally within a given discourse. The third task, idiom cloze, challenges models to retrieve an idiom suitable for a specific situation to fill a masked position in a paragraph. Together, these tasks are designed to capture real-world idiom understanding, emphasizing both contextual reasoning and functional usage.

\subsection{The Challenge of Persian Idioms}
Persian idioms pose specific, non-trivial challenges for NLP systems that distinguish them from idioms in English or other European languages.
\begin{itemize}
    \item \textbf{Morphological Diversity:} Persian maintains a complex inflectional morphology. A verbal idiom can appear in various tenses, aspects, and moods. Furthermore, Persian syntax allows for significant scrambling; the components of an idiom may be separated by other words (e.g., objects, adverbs), making exact string matching insufficient for detection.
    \item \textbf{Diglossia:} Persian functions in a diglossic environment where the formal written register (\textit{ketabi}) and the colloquial spoken register (\textit{mohaverei}) differ significantly. Idioms often bridge these registers, appearing in formal news text with colloquial morphology, or vice versa, requiring models to be robust to stylistic variation.
    \item \textbf{Evaluative Connotation:} A significant proportion of Persian idioms are ``evaluative'' in nature; they are used to judge behavior, character, or situations. As annotation reveals, this evaluation is predominantly negative; idioms serve as tools for social regulation, warning, or criticism. This creates a label imbalance that models must learn to navigate.
\end{itemize}

\subsection{Contributions}
PersianIdioBench addresses these critical gaps by moving beyond simple lexicon curation to establish a comprehensive evaluation framework. This paper makes the following primary contributions:
\begin{itemize}
    \item \textbf{Resource Construction:} The largest dataset of Persian idioms to date is compiled and refined, comprising 4,352 unique entries, by crawling web dictionaries such as Abadis, encyclopedic sources including Dehkhoda and Amid, parsing online educational materials (for example Paarsadab and Dyanati), and consolidating previously published academic lists.
    \item \textbf{Gold-Standard Annotation:} A high-quality, human-verified layer of evaluative connotation (sentiment) is provided for these idioms, achieving substantial inter-annotator agreement ($\kappa=0.68$).
    \item \textbf{Contextual Benchmark Design:} A robust set of context-aware tasks is created, mining thousands of usage examples from diverse corpora (news, social media, forums) to support Appropriateness and Cloze tasks.
    \item \textbf{Comprehensive Evaluation:} A range of LLMs is benchmarked, identifying a ``pragmatic gap'' where models fail to detect subtle misuse of idioms, establishing a baseline for future research in Persian NLP.

The resulting benchmark and datasets are publicly released to serve as a standardized evaluation resource.
\end{itemize}

\section{Related Work}

\subsection{LLM Evaluation}
The development of PersianIdioBench is positioned at the intersection of LLM evaluation, figurative language processing, and research on Persian as a low-resource language. The emergence of diverse benchmarks has substantially advanced the systematic evaluation of large language models. Early efforts such as GLUE \cite{wang2018glue} and SuperGLUE \cite{wang2019superglue} established foundational standards for measuring language understanding. Subsequent benchmarks including MMLU \cite{hendrycks2020mmlu}, HELM \cite{liang2022holistic}, BigBench \cite{srivastava2022beyond}, HellaSwag \cite{zellers2019hellaswag}, ARC \cite{clark2018arc}, and MMLU-Pro \cite{wang2024mmlupro} expanded evaluation coverage. More recently, culturally grounded benchmarks have been introduced to assess LLM performance in underrepresented languages, including Arabic \cite{qian2024cameleval, nacar2025inclusive}, Korean \cite{kim2024click}, and Russian \cite{vasilev2025ruscode}. In the context of Persian, steady progress has been made in developing resources for both model training \cite{hosseinbeigi2025matina, sabouri2022naab} and evaluation \cite{hosseinbeigi2025persianeval, shamsfard2025farseval}. Additionally, \citet{ghahroodi2024persianmmlu} presented a Persian adaptation of the MMLU benchmark. This work continues this line of research by introducing a benchmark specifically designed to evaluate LLMs on Persian idiomatic language.



\subsection{Figurative Language Benchmarks}
In contrast to psycholinguistic research, computational linguistics studies aim to model idioms in text through tasks including detection, disambiguation, and paraphrasing \cite{flor2025idiomsurvey}. A substantial body of work focuses on binary classification tasks that distinguish figurative from literal usages in context, such as the VNC-Tokens dataset \cite{fazly2009idiomidentification}, the MAGPIE corpus \cite{haagsma2020magpie}, the RU Idioms dataset \cite{aharodnik2018russianidiom}, and the DICE dataset \cite{mi2025rollingdice}. Research has also extended idiom processing to affective meaning via IDIOMENT \cite{williams2015idiomsentiment}, SLIDE \cite{jochim2018slide}, and IDEM \cite{prochnow2024idem}. Building on idiom understanding, several studies have introduced datasets that support multilingual evaluation, such as MABL \cite{kabra2023multicultural}, MAPS \cite{liu2023culturalproverbs}, and IdiomKB \cite{li2024idiomkb}. Beyond classification and sentiment tasks, a growing line of work targets generation-oriented capabilities. The Parallel Idioms Corpus \cite{zhou2021pie} and the Idiom Substitution dataset \cite{liu2016phrasalsub} support modeling of semantic equivalence. Complementary benchmarks further evaluate productive usage: the ChID dataset \cite{zheng2019chid} and the CIP dataset \cite{qiang2023idiomparaphrase}. More recently, CHENGYU-BENCH \cite{fu2025chengyubench} introduced a multi-task evaluation suite that includes evaluative connotation classification, contextual appropriateness detection, and an open cloze task without answer options. CHENGYU-BENCH serves as a key methodological inspiration for this work, although Persian's morphological diversity introduces distinct challenges not present in fixed four-character Chinese expressions.

\begin{table}[t]
\centering
\caption{Source Distribution of the Final Repository}
\label{tab:source_distribution}
\begin{tabular}{lcc}
\toprule
\textbf{Source} & \textbf{Count} & \textbf{Percentage} \\
\midrule
Abadis Proverbs & 2,619 & 60.2\% \\
PersianIdioms & 1,112 & 25.6\% \\
Moein Idioms & 312 & 7.2\% \\
Paarsadab & 256 & 5.9\% \\
Ostovar & 51 & 1.2\% \\
Dyanati & 2 & 0.05\% \\
\bottomrule
\end{tabular}
\end{table}

\subsection{Persian NLP Resources}
Historically, there have been few NLP resources for Persian in comparison to other languages. Traditional lexicography provided the foundation, but dictionaries like Dehkhoda and Moein lack the structured format required for machine learning. Computational efforts include:
\begin{itemize}
    \item \citet{ostovar2014persianidioms}: Compiled a database of 364 verb-centric idioms, pivotal for translation studies but too limited in scale for modern LLM evaluation.
    \item \citet{dashtipour2022extending}: Extended the PerSent lexicon with $\sim$1,000 idioms to improve sentiment analysis, demonstrating that ignoring idioms harms performance.
    \item \citet{rezaeimanesh2025largelanguagemodelspersian}: Introduced \textit{PersianIdioms}, a lexicon of $\sim$2,300 idioms with definitions, evaluated primarily in the context of machine translation.
        \item \citet{khoshtab2025multilingualidioms}: Introduced Persian evaluation sets for idiom and simile interpretation within a broader multilingual study.
    \item \citet{farsi2025melac}: Introduced a Persian idiom evaluation task, Proverbs-Quiz, consisting of 370 samples where models must select the correct meaning of a proverb.
\end{itemize}
PersianIdioBench integrates these prior resources into a larger, unified framework, adding thousands of new entries from web crawling and providing the first dedicated benchmarks for appropriateness and cloze generation in Persian.

\section{Datasets}
The foundation of the benchmark is a comprehensive, deduplicated repository of Persian idioms.

\subsection{Data Acquisition Sources}
Four primary categories of sources were targeted to build the initial pool of expressions.

\subsubsection{Abadis Web Dictionary}
The Abadis website \cite{abadi_dictionary} serves as a central hub for Persian lexicographic resources and explicitly permits the use of its data provided that proper attribution is given. Leveraging this policy, a stateful web crawler was developed to systematically traverse all proverb and idiom pages and ensure complete coverage of the content made available on the platform. Through this process, 2,690 idioms and proverbs were extracted, which constitute the backbone of the initial collection due to the richness and reliability of their metadata. For each entry, Abadis provides detailed information including the idiom’s meaning, user-contributed comments, and, when available, references to credible dictionaries such as Moeen, Dehkhoda, and Amid.

\subsubsection{Educational Resources (Paarsadab and Dyanati)}
To capture idioms frequently taught in the Persian education system, two major educational repositories were identified: Paarsadab \cite{paarsadab_dictionary} and Dyanati \cite{dyanati_resources}. The available PDF documents were converted into machine-readable text and structured tables were extracted, yielding 756 proverbs from Paarsadab and 375 proverbs from Dyanati. In addition, several other PDF resources were identified, including books and compilations such as encyclopedic collections of idioms and proverbs. However, these materials require sophisticated OCR pipelines due to complex layouts and scan quality. Their integration is left as future work to further expand and enrich the dataset.

\subsubsection{The Moein Dictionary}
The Moein Dictionary is one of the most reputable lexicographic resources for Persian and provides extensive lexical and semantic descriptions. Moein entries were accessed through their structured representation on the Abadis platform, which offers detailed metadata similar to the proverb entries described earlier, including meanings, explanations, and cross-references to other dictionaries. Following the linguistic assumption that idioms are inherently multi-word expressions (MWEs) \cite{findlay-2017-multiword}, entries consisting of more than one token were extracted, resulting in 5,105 candidate expressions. These MWEs were treated as potential idioms to be incorporated into the resource pool. To further filter for true idiomatic usages, a scoring heuristic was implemented: a score of (+2) was assigned to entries whose definitions contained indicative keywords such as kenayeh (metaphor) or estelah (idiom). Conversely, entries were penalized by (-1) if their source indicated non-figurative domains (e.g., the Academy of Persian Language and Literature or encyclopedic name lists). Retaining only entries with a positive aggregate score yielded 1,183 high-probability idioms.

\subsubsection{Existing Academic Datasets}
The \textit{PersianIdioms} dataset (2,334 entries) and the \citet{ostovar2014persianidioms} collection (364 entries) were incorporated to ensure backward compatibility and completeness.

\subsection{Normalization and Deduplication Pipeline}
Merging these heterogeneous sources required a rigorous normalization and deduplication pipeline. Text normalization was first applied by removing digits and punctuation, unifying character variants, and eliminating artifacts. The text was tokenized and lemmatized using the Hazm library \cite{roshan2014hazm}. For tokenization consistency, the ParsBERT tokenizer \cite{farahani2021parsbert} was additionally applied to the normalized text. Pairwise duplicate detection was performed using morphology-oriented similarity measures (Jaccard similarity, Levenshtein distance, SequenceMatcher similarity, and trigram cosine similarity). Pairs exceeding thresholds were flagged, and a canonical instance was retained according to a predefined source-reliability hierarchy.

\subsection{Final Dataset Statistics}
The deduplication process identified 4,299 duplicate pairs across sources, along with 123 hard duplicates occurring within individual datasets. After merging and cleaning, the final \textbf{PersianIdioBench} repository contains \textbf{4,352 unique idioms}. Table~\ref{tab:source_distribution} presents the distribution of entries across sources.

\section{Benchmark}
\subsection{Task Definition}
To comprehensively assess a model’s ability to understand and use idiomatic expressions in realistic contexts, this benchmark is composed of three complementary subtasks:

\textbf{Evaluative Connotation.} Accurately identifying an idiom’s intrinsic polarity, positive or negative, is essential for understanding idiomatic language. In this subtask, models are required to predict the sentiment polarity conveyed by each idiom in isolation (Table~\ref{tab:idiom_examples}).

\textbf{Appropriateness.} Determining whether a Persian idiom is used correctly in context requires evaluating its semantic compatibility, sentiment orientation, and pragmatic force. This subtask evaluates whether a model can judge the contextual suitability of an idiom within a paragraph.

\textbf{Cloze.} This subtask assesses a model’s ability to select an idiom that best completes a paragraph-level context. For each instance, a span in the text is masked, and the model is presented with 10 options: one correct idiom and nine distractors. The model is asked to rank these options.

\begin{table*}[t]
\centering
\caption{Surface meaning, figurative meaning, and polarity of example Persian idioms}
\label{tab:idiom_examples}
\begin{tabular}{p{3cm}p{3.2cm}p{5.2cm}c}
\toprule
\textbf{Idiom} & \textbf{Surface Meaning} & \textbf{True Meaning} & \textbf{Polarity} \\
\midrule
gol be sabzeh zadan & Hitting a flower onto greenery & Doing something perfectly or adding an elegant finishing touch & Positive \\
\addlinespace
namak nashenas & Someone who does not recognize salt & An ungrateful person who does not appreciate kindness & Negative \\
\addlinespace
sang jelo pash andakhtan & Throwing a stone in front of someone’s feet & Creating obstacles or deliberately hindering someone’s progress & Negative \\
\bottomrule
\end{tabular}
\end{table*}

\subsection{Connotation Annotation}
Two native Persian speakers were recruited to annotate the full set of 4,352 idioms into four categories: Positive, Negative, Neutral, and Not idiom. Neutral and Not idiom classes were ultimately merged. The inter-annotator agreement achieved was substantial (Cohen's $\kappa = 0.68$). Conflicting labels were resolved through consensus discussions between the annotators and an independent third-party expert. This process yielded a heavily negative-skewed distribution (52.1\% Negative, 13.1\% Positive), consistent with the nature of Persian social commentary. For benchmark experiments, evaluation is performed on the binary distinction between positive and negative idioms, resulting in 2,839 evaluation pairs.

\subsection{Context Mining and Corpus Compilation}
To construct datasets for the Appropriateness and Cloze tasks, a large corpus of Persian text was compiled from diverse domains, including news and social media. The source datasets comprise Hamshahri \cite{aleahmad2009hamshahri}, Tasnim \cite{pourmand2022tasnim}, Asriran \cite{pourmand2022asriran}, Tarjoman \cite{pourmand_tarjoman}, EmoPars \cite{sabri2021emopars}, ParsABSA \cite{ataei2019parsabsa}, DeepSentiPers \cite{sharami2020deepsentipers}, and ArmanEmo \cite{mirzaee2025armanemo}, collectively providing broad linguistic coverage across both formal and informal registers.

Since idioms often appear with minor lexical, morphological, or syntactic variations, fuzzy matching rather than exact string search was employed to retrieve idioms from text. Specifically, the \texttt{RapidFuzz} library \cite{rapidfuzz_token_set_ratio} was used, which provides efficient approximate string-matching functions. Two complementary similarity metrics were applied: \texttt{token\_set\_ratio}, which compares unique words between the idiom and text regardless of order to capture reordered or slightly modified expressions, and \texttt{partial\_ratio}, which searches for the optimal alignment of the idiom within longer passages and computes an edit-distance-based similarity score.

Scores range from 0 to 100, and a text was considered a match if both metrics exceeded 85, or if \texttt{partial\_ratio} reached 100, indicating a near-exact local match. The metric threshold was set to 85 to balance recall and precision, accommodating morphological variations without introducing excessive false positives. All comparisons were parallelized to improve processing efficiency. For long documents (Hamshahri, Tasnim, Asriran, Tarjoman), up to seven paragraphs surrounding each idiom occurrence were extracted (three preceding and three following) to preserve sufficient context while minimizing irrelevant content.

After initial retrieval, deduplication and merging were performed, resulting in 1,012 unique idioms appearing in 5,923 candidate contexts. To ensure that matches reflected genuine idiomatic usage rather than incidental word overlap, an automatic verification step was applied using GPT-4o-mini, prompting the model to confirm whether the idiom functioned as a semantic unit within the text. This reduced the dataset to 3,446 validated text–idiom pairs. Contexts shorter than ten words were subsequently removed to ensure sufficient discourse information for pragmatic evaluation. The final corpus contains 3,425 contexts covering 839 unique idioms, exhibiting a long-tail frequency distribution where most idioms appear in few contexts. 
Table~\ref{tab:context_pairs} summarizes the number of text–idiom matches obtained from each dataset, highlighting the corpus coverage achieved through the fuzzy matching and validation pipeline.

\begin{table}[htbp]
\centering
\caption{Number of extracted text–idiom pairs by dataset}
\label{tab:context_pairs}
\begin{tabular}{lc}
\toprule
\textbf{Dataset} & \textbf{Pairs} \\
\midrule
Asriran & 862 \\
Hamshahri & 362 \\
Tasnim & 508 \\
Tarjoman & 361 \\
EmoPars & 226 \\
ParsABSA & 81 \\
DeepSentiPers & 17 \\
ArmanEmo & 89 \\
\midrule
\textbf{Total} & 5,923 \\
\bottomrule
\end{tabular}
\end{table}

\subsection{Task 2: Appropriateness in Context}
This task evaluates whether a model can distinguish between appropriate and inappropriate usage of an idiom within a given context. The dataset is constructed as a binary classification problem using both naturally occurring and adversarially generated examples.

\textbf{Positive Samples.}
Positive instances are drawn directly from the validated context–idiom pairs obtained through the context mining pipeline. Each text contains an idiom appearing in a coherent and meaningful context, regardless of whether the usage is literal or figurative. These instances therefore represent naturally occurring appropriate usages.

\textbf{Negative Samples.} Negative instances were generated by replacing the original idiom in each context with a semantically incompatible distractor idiom. Candidate distractors were identified using fuzzy string matching with the \texttt{token\_sort\_ratio} from the RapidFuzz library \cite{rapidfuzz_token_set_ratio}, ensuring a high degree of lexical overlap to prevent models from solving the task through simple surface-level detection. To ensure the challenge remained pragmatic rather than purely lexical, a polarity-mismatch constraint was enforced: selected distractors had to possess an evaluative connotation opposite to the original idiom. This ensures that while the distractor may look similar to the original (e.g., sharing keywords), it creates a clear semantic and pragmatic clash with the surrounding discourse, forcing the model to rely on deep contextual reasoning rather than keyword matching.

The final Appropriateness dataset contains 6,057 instances and is balanced between appropriate and inappropriate usages, ensuring that performance reflects models’ ability to evaluate contextual fit rather than label frequency. Table~\ref{tab:app_samples} provides representative examples.

\begin{table*}[htbp]
\centering
\caption{Example instances from the Appropriateness dataset with idioms in context.}
\label{tab:app_samples}
\begin{tabular}{p{3cm} p{3cm} p{7cm} c}
\toprule
\textbf{Idiom in Text} & \textbf{Target Idiom} & \textbf{Context (English translation with Pinglish idiom)} & \textbf{Label} \\
\midrule
cheshm zagh \newline(light-colored eyes) & cheshm o cheragh \newline(beloved, the pride of) & 
When I wrote the script and told him he should play the lead, he laughed and said: I neither have blond hair nor **cheshm zagh**. How do you want a worn-out peach to play the main role? & Inapp. \\
\addlinespace
cheshm o cheragh \newline(beloved, the pride of) & cheshm o cheragh \newline(beloved, the pride of) &
We repeatedly wrote letters and asked why newspapers did not investigate what has been done for the families of martyrs, the **cheshm o cheragh** of the nation. & App. \\
\bottomrule
\end{tabular}
\end{table*}

\subsection{Task 3: Cloze-Style Idiom Selection}
This task evaluates a model’s ability to identify the idiom that best fits a paragraph-level context. Each instance consists of a naturally occurring text in which the target idiom is replaced by a placeholder token (\texttt{<<MASK\_TOKEN>>}), and the model must rank a fixed set of candidate idioms according to their contextual suitability.

\textbf{Context Preparation.}
This dataset was constructed from the subset of mined contexts in which the idiom appeared as a contiguous string exactly matching its canonical form. Restricting to exact matches simplifies masking and avoids ambiguity introduced by discontinuous or morphologically altered occurrences. This filtering step resulted in 2,457 eligible texts. In each text, the first occurrence of the idiom was replaced with the mask token.

\textbf{Candidate Set Construction.}
For each paragraph–idiom pair, 50 idioms were first sampled from the full idiom inventory to create a candidate pool. This pool, along with the true idiom, was then provided to the strongest available GPT-5 series model at the time (GPT-5.2). The model was instructed, given knowledge of the correct answer, to select 10 idioms that would be irrelevant or misleading in the given context while remaining distinct from the gold idiom. To mitigate potential ``model-in-the-loop'' bias, where OpenAI models might find OpenAI-generated distractors more predictable, all candidate sets underwent a rigorous manual review by native speakers. Annotators refined the sets to ensure distractors were not only contextually inappropriate but also linguistically diverse, removing any model-specific stylistic artifacts or repetitive phrasing that could serve as ``shortcuts" for specific model architectures. Models are required to rank all candidates from most to least suitable for the masked position. The ranking formulation provides a fine-grained evaluation signal by capturing relative preference ordering among competing idiomatic expressions. The final Cloze dataset contains 2,457 instances.

\section{Experimental Setup}

\subsection{Models Evaluated}
A set of models representing both modern closed-source LLMs and widely used open-source Persian language models is evaluated.

\textbf{Closed-Source Models.}
GPT-5-mini and Claude-Haiku-3.5 are evaluated under prompt-based settings. Both models are tested in zero-shot configurations across all tasks, and in few-shot settings for the binary classification tasks (Connotation and Appropriateness). Larger proprietary models are intentionally excluded due to their substantially higher cost, as the primary goal is comparative benchmarking rather than exhaustive leaderboard optimization.

\textbf{Open-Source Models.}
Two representative Persian-focused baselines are included:

\begin{itemize}
    \item Dorna \cite{dorna_llama3}. A LLaMA-based Persian language model. It is evaluated under zero-shot, few-shot, and chain-of-thought (CoT) prompting. Additionally, lightweight supervised instruction fine-tuning using QLoRA \cite{dettmers2023qlora} is performed.
    \item ParsBERT \cite{farahani2021parsbert}. A BERT-based Persian encoder model. For classification tasks, the pretrained encoder is frozen and a linear classification head is trained. For the appropriateness task, a cross-encoder formulation is adopted by concatenating the idiom and paragraph as a single sequence.
\end{itemize}

ParsBERT is not evaluated on the Cloze task because its masked-token prediction objective is not directly compatible with ranking multi-token idiom candidates without additional architectural modifications.

\subsection{Data Splits}
Across all tasks, a small portion of the dataset (10\%) is reserved for model adjustment purposes such as few-shot example selection or fine-tuning, while the remaining 90\% is used as the held-out evaluation benchmark. This split deviates from conventional training-heavy setups because the primary goal is to establish a reliable benchmark rather than maximize task-specific performance through extensive training. The held-out portion remains untouched to ensure fair and unbiased comparison across models.

\subsection{Task-Specific Settings}
For the connotation task, it is framed as binary sentiment classification (positive vs.\ negative). Closed-source models are evaluated in zero-shot and few-shot settings, while Dorna is additionally evaluated with chain-of-thought (CoT) prompting and fine-tuning to exploit reasoning. For ParsBERT, a cross-encoder architecture is used where the input is structured as \texttt{[CLS] idiom [SEP] paragraph [SEP]}. The [CLS] token allows aggregation of the sequence-level representation, while the [SEP] tokens explicitly separate the idiom from its surrounding context, helping the classifier focus on the relevant interaction.

In the appropriateness task, models predict whether an idiom is correctly used in a paragraph. For ParsBERT, the frozen encoder is paired with a linear classification head, ensuring that only task-specific adaptation occurs while keeping the pretrained language representation stable. For Dorna, a causal language modeling formulation is adopted rather than adding an MLP classification head on top of the Persian LLaMA backbone. This choice is motivated by the small size of the dataset and the difficulty of the task; training a new head from scratch would be unstable and less effective. The model generates token probabilities in sequence, allowing the attention layers to leverage contextual cues for evaluating appropriateness.

For the cloze task, each paragraph contains a masked idiom (\texttt{<<MASK\_TOKEN>>}) and the model ranks candidate idioms. Closed-source models are evaluated in zero-shot settings, while Dorna is tested in zero-shot, CoT, and fine-tuned settings. The retrieval formulation allows direct comparison of the predicted idioms to the gold idiom without requiring task-specific heads. Few-shot prompting is not used due to instability in ranking outputs. All these architectural choices, including causal LM for Dorna and frozen encoder for ParsBERT, are motivated by limited data and the need to leverage pretrained knowledge efficiently.

\subsection{Training Configuration}
Dorna is fine-tuned using supervised instruction tuning with QLoRA. Only the query and value projection matrices (\texttt{q\_proj}, \texttt{v\_proj}) of the attention layers are updated, which reduces computational overhead while allowing effective task adaptation. To avoid overfitting on the small dataset, a low-rank LoRA setup is employed with a rank ($r$) of 4, an alpha of 8, and a dropout rate of 0.15. Training is performed on a Tesla T4 GPU for 1 epoch using a batch size of 4, a learning rate of $2 \times 10^{-5}$, and a maximum sequence length of 256. The task is formulated as causal language modeling (Causal LM).

For ParsBERT, the pretrained encoder is frozen and only a linear classification head is trained for the connotation and appropriateness tasks. While full fine-tuning is common, preserving the stable, broad-spectrum language representations of the original model was prioritized given the targeted and nuanced nature of idiomatic processing. This approach prevents the model from overfitting on the limited 10\% tuning split and ensures that the reported performance reflects the model's inherent ``idiomatic knowledge'' rather than a localized adaptation to the benchmark's specific vocabulary. Fine-tuning is conducted for 1 epoch utilizing a batch size of 16, a learning rate of $1 \times 10^{-5}$, a weight decay of 0.01, and a warmup ratio of 0.1. The maximum sequence length is set to 256 and bf16 precision is enabled for mixed-precision training.

\subsection{Evaluation Metrics}
For Connotation and Appropriateness, Accuracy, Macro F1, Matthews Correlation Coefficient (MCC), and Cohen’s Kappa ($\kappa$) are reported. For Cloze, Top-1 Accuracy and Mean Reciprocal Rank (MRR) are reported.

\section{Results and Analysis}

\textbf{Task 1: Connotation Classification Results.}
Table~\ref{tab:connotation_results} shows that OpenAI’s model achieves the strongest overall performance. Claude benefits noticeably from few-shot examples. Open models remain clearly below the closed models, and ParsBERT operates below the random baseline, implying frozen representations with a simple classifier are insufficient for this task.

\begin{table}[t]
\centering
\small
\caption{Connotation Classification Performance (Macro F1)}
\label{tab:connotation_results}
\begin{tabular}{llcc}
\toprule
\textbf{Model} & \textbf{Setting} & \textbf{Macro F1} & \textbf{MCC} \\
\midrule
OpenAI & Zero-Shot & 0.776 & 0.555 \\
OpenAI & Few-Shot & \textbf{0.780} & \textbf{0.562} \\
Claude & Zero-Shot & 0.661 & 0.404 \\
Claude & Few-Shot & 0.762 & 0.525 \\
Dorna & Zero-Shot & 0.542 & 0.206 \\
Dorna & Few-Shot & 0.603 & 0.253 \\
Dorna & Fine-Tuned & 0.521 & 0.265 \\
ParsBERT & Fine-Tuned & 0.464 & 0.039 \\
\midrule
Random Baseline & - & 0.500 & 0.000 \\
\bottomrule
\end{tabular}
\end{table}

\textbf{Task 2: Appropriateness Results.}
Table~\ref{tab:appropriateness_results} indicates a clear gap between OpenAI’s GPT-5-mini and other models. Claude correctly identifies appropriate usages far more often than inappropriate ones, indicating that pragmatic violations remain difficult to detect. ParsBERT and Dorna operate close to the random baseline.

\textbf{Task 3: Cloze Ranking Results.}
Table~\ref{tab:cloze_results} shows OpenAI’s model achieves the strongest performance on the ranking task. Claude performs noticeably lower in zero-shot. Dorna exhibits modest zero-shot results, with fine-tuning yielding only a small improvement.

\section{Analysis \& Discussion}

\subsection{Qualitative Analysis}
To quantify instance difficulty, an agreement-based score, $n_{\text{correct}} \in \{0, \dots, 3\}$, is defined, representing the number of models (OpenAI, Claude, Dorna) that produce the correct prediction. Representative easy and hard cases are analyzed to identify systematic error patterns.

\textbf{Appropriateness Judgments.} 
Models succeed when clear semantic or register cues reinforce the idiom. For example, \textit{ghor zadan} (``to complain'') and \textit{cheragh sabz neshan dadan} (``to give the green light'') are correctly identified when anchored by specific lexical cues (e.g., children's behavior or institutional authorization). Conversely, hard cases involve subtle pragmatic mismatches. Models struggle when idioms share surface imagery (e.g., \textit{tooti meqal} vs.\ \textit{tooti-vaar}) or when abstract contexts provide weak explicit cues (e.g., \textit{dari be takhte khordan}). Errors in \textit{bal-o-par dadan} vs.\ \textit{az chaleh be chah oftadan} further suggest that models fail when an incorrect idiom remains event-compatible but alters agency relations.

\textbf{Cloze Task.} 
Easy instances follow a similar pattern where strong syntactic or lexical constraints license the idiom. For instance, \textit{gooshesh bedehkar nist} (``indifferent'') is successfully predicted when contexts describe ignored warnings. Harder cases reveal errors driven by \textbf{lexical overlap} and \textbf{sentiment misalignment}. Models frequently select \textit{cheshm zakhm} (``evil eye'') over \textit{cheshm-e zagh} (``light eyes'') due to the shared token \textit{cheshm} (eye). Additionally, models often select negative idioms (e.g., \textit{khak khordan}) for \textit{cheshm-o-del sir} (``satiated'') because they match the general sentiment but fail to capture the specific experiential meaning.

\subsection{The Pragmatic Gap Revisited}
Beyond aggregate scores, the results reveal a consistent discrepancy between semantic recognition and pragmatic reasoning. Although models perform strongly on connotation classification, where sentiment cues are often lexically grounded, performance declines on appropriateness and Cloze selection tasks. Both require the integration of discourse context, implicit causality, and register constraints.

This gap suggests that idiom understanding is not a single capability but a layered process. Models frequently retrieve the coarse-grained semantic polarity of an expression yet struggle when correct usage depends on subtle contextual alignment. Errors often occur in cases where an idiom remains event-compatible but introduces incorrect agency, intent, or discourse function. Such patterns indicate a reliance on surface semantic similarity rather than robust pragmatic modeling. Taken together, these findings reinforce the view that pragmatic competence remains a primary bottleneck for LLMs in figurative language. This is particularly evident in low-resource settings, where exposure to diverse conversational contexts is limited, hindering the model's ability to move beyond literal or high-level semantic approximations.

\begin{table}[htpb]
\centering
\small
\caption{Cloze Task Ranking Metrics}
\label{tab:cloze_results}
\begin{tabular}{llcc}
\toprule
\textbf{Model} & \textbf{Setting} & \textbf{Top-1 Acc} & \textbf{MRR} \\
\midrule
OpenAI & Zero-Shot & \textbf {69.2\%}& \textbf {0.785} \\
Claude & Zero-Shot & 45.8\% & 0.594 \\
Dorna & Zero-Shot & 21.1\% & 0.318 \\
Dorna & Fine-Tuned & 22.2\% & 0.329 \\
\midrule
Random Baseline & - & 9.1\% & 0.275 \\
\bottomrule
\end{tabular}
\end{table}

\section{Limitations}
This benchmark has several constraints that should be considered. First, rare dialectal or highly domain-specific idioms may be underrepresented, leading to a coverage bias that favors more common expressions. Additionally, the adversarial examples generated for the appropriateness task may not perfectly reflect the natural distribution of errors occurring in authentic discourse. Finally, the benchmark’s scope primarily evaluates discriminative and retrieval-style competence, and does not fully assess open-ended generation capabilities.

\begin{table}[htbp]
\centering
\small
\caption{Appropriateness Task Performance}
\label{tab:appropriateness_results}
\begin{tabular}{llcc}
\toprule
\textbf{Model} & \textbf{Setting} & \textbf{Macro F1} & \textbf{Cohen's $\kappa$} \\
\midrule
OpenAI & Zero-Shot & \textbf{0.828} & \textbf{0.657}\\
OpenAI & Few-Shot & 0.815 & 0.630 \\
Claude & Zero-Shot & 0.646 & 0.311 \\
Claude & Few-Shot & 0.658 & 0.343 \\
Dorna & Zero-Shot & 0.496 & 0.007 \\
ParsBERT & Cross-Encoder & 0.473 & 0.004 \\
\bottomrule
\end{tabular}
\end{table}

\section{Conclusion}
PersianIdioBench provides a comprehensive multi-task benchmark designed to evaluate large language models on Persian idiom understanding across semantic, pragmatic, and contextual dimensions. By combining a large curated idiom repository with high-quality annotations, a structured framework is provided for analyzing figurative language competence in a low-resource setting. The experiments reveal that while advanced models demonstrate strong zero-shot capabilities, all systems exhibit a measurable decline as tasks demand deeper pragmatic reasoning, highlighting the persistent gap between lexical knowledge and discourse-level understanding.

\section*{Author Contributions}
\begin{itemize}
  \item Reyhaneh Hashempour --- Conceptualization; Data curation; Formal analysis; Investigation; Methodology; Resources; Software; Validation; Visualization; Writing -- original draft; Writing -- review \& editing.
  \item Richard Sutcliffe --- Methodology; Supervision; Writing -- review \& editing.
  \item Jon Chamberlain --- Methodology; Supervision; Writing -- review \& editing.
\end{itemize}

\bibliographystyle{acl_natbib}
\bibliography{ref}

\end{document}